This study is limited to an analysis of data within the biomedical field available through PubMed, and it remains unclear whether similar results and discussions would be obtained when applied to other or broader disciplines.

As one validation, we compared the citation counts for the two most-cited papers mentioned earlier between this study and "The top 100 papers" covering the entire natural sciences (comparing this study's 1946--2015 cumulative generations with Web of Science's 2014 data). For "Protein measurement with the Folin phenol regent," the counts were 49,164 and 305,148, respectively, and for "Cleavage of structural protein during the assembly of the head of bacteriophage T4," they were 49,916 and 213,005, respectively. This represents a difference of approximately four to six times, indicating that direct interpretation is not appropriate. However, a comparison of the paper titles present in both datasets revealed a substantial number of matches.

From a journal-level perspective, "The top 100 papers" also demonstrates that high-Impact Factor (IF) journals do not necessarily dominate the list. In this regard, it is reasonable to assume that interdisciplinary fields share similar characteristics with the biomedical domain.

Regarding theme analysis, from the perspective of adapting this method to broader disciplines, it is difficult to utilize controlled vocabularies as systematically organized as MeSH descriptors in other databases, making analysis with equivalent precision to this study challenging.

